% ECONOMICS_PROBLEM: {"id": "G12-324", "title": "Intermediary constraints and risk pricing", "jel_code": "G12", "source_id": "OP-0316"}
% !TeX program = xelatex
\documentclass[11pt,a4paper]{article}
\usepackage[margin=23mm]{geometry}
\usepackage{fontspec}
\setmainfont{Latin Modern Roman}
\usepackage{amsmath,amssymb,mathtools}
\usepackage{xcolor,enumitem,booktabs,longtable,array,xurl,hyperref}
\definecolor{muted}{HTML}{53636C}
\hypersetup{colorlinks=true,urlcolor=blue,linkcolor=blue}
\setlength{\parindent}{0pt}
\setlength{\parskip}{5pt}
\newcommand{\R}{\mathbb R}
\newcommand{\E}{\mathbb E}
\newcommand{\N}{\mathbb N}
\newcommand{\ind}{\mathbf 1}
\newcommand{\Var}{\operatorname{Var}}
\newcommand{\Cov}{\operatorname{Cov}}
\newcommand{\supp}{\operatorname{supp}}
\DeclareMathOperator*{\argmin}{arg\,min}
\DeclareMathOperator*{\argmax}{arg\,max}
\newcommand{\entrymeta}[3]{{\sffamily\small\color{muted}\textbf{\texttt{#1}}\quad|\quad #2\par #3\par}}
\newcommand{\Problemref}[1]{\texttt{#1}}
\newcommand{\JELref}[2]{\texttt{#2} (JEL \texttt{#1})}
\begin{document}
\section*{Intermediary constraints and risk pricing}
\textbf{Problem ID:} \texttt{G12-324}\quad \textbf{JEL:} \texttt{G12}\par
% BEGIN ECONOMICS STATEMENT
\label{problem:OP-0316}
{\sffamily\small\color{muted}Original subject group: Asset pricing and financial markets\par}
\label{definitions:problem:30007598}
\textbf{Audit classification:} Published research agenda. The official agenda supports financial-actor/regulatory-constraint research broadly. The funding instrument, multiplier identification and particular asset-pricing experiment are editorial.
\subsection*{Definitions and precise research target}
Consider assets with real gross return vector \(R_{t+1}\), households, and intermediaries subject to balance-sheet constraints. Intermediary \(i\) has equity \(e_{it}\), leverage bound \(\ell_{it}\), and a Lagrange multiplier \(\lambda_{it}\geq0\) on its financing constraint. Let \(z_t\) be an exogenous shock to intermediary funding capacity that does not directly change asset cash flows or household preferences. Define state-dependent causal responses
\[\tau_j(h,x)=\frac{\partial E[R^j_{t+h}(do(z_t=u))\mid x]}{\partial u}\Big|_{u=0},\]
where \(h\geq1\) is a horizon and \(x\) records initial constraints and market segmentation. The task is to identify these responses and characterize when intermediary marginal utilities or constraint multipliers determine equilibrium risk prices. Establish restrictions linking observed leverage, equity and funding spreads to latent \(\lambda_{it}\), including the possibility that some intermediaries are unconstrained. State why the funding instrument satisfies exclusion, and account for endogenous entry and household substitution. Compare normal and crisis states without imposing linearity at constraint boundaries. Return an identified set when different pricing kernels or cash-flow channels fit the same data. Covariation between intermediary leverage and returns alone does not establish the marginal-pricing mechanism.

\textbf{Additional formal definitions and scope (editorial).} Take \(e_{it}>0\), asset value \(A_{it}\geq0\), and a leverage restriction \(A_{it}\leq\ell_{it}e_{it}\), where \(\ell_{it}\geq1\). Its multiplier satisfies \(\lambda_{it}\geq0\), slackness \(\lambda_{it}(\ell_{it}e_{it}-A_{it})=0\), and the model's portfolio first-order conditions. Specify the objective and financing prices before interpreting \(\lambda\) as a shadow value; it is not an observed funding spread by definition. The shock intervention alters a stated limit, equity endowment, or funding supply equation; these are different shocks. Outcome \(R^j_{t+h}\) is the gross return over dates \(t+h-1,t+h\); also distinguish its predictable risk premium from its unexpected repricing. At a binding boundary use one-sided responses or finite differences if derivatives fail to exist.

For the identification part, declare an admissible class \(\Theta\) of structural models, including preferences, technologies, shock laws, measurement/sampling rules and any equilibrium selector. If \(O\) denotes the complete observable history and \(T(\theta)\) the target defined above, the identified set is \(\mathcal I_T=\{T(\theta):\theta\in\Theta,\ \mathcal L_\theta(O)=\mathcal L_{\rm obs}(O)\}\); point identification means this set is a singleton. A \(do(a)\) comparison replaces stated policy/technology equations while fixing the declared exogenous law and re-solving the remaining equations. These definitions do not assert that an identification theorem holds without further restrictions.
\subsection*{Variants and assumption changes}
\begin{enumerate}
\item \textbf{Slack constraint (editorial).} Restrict to states with \(A_i<\ell_i e_i\) and \(\lambda_i=0\); distinguish funding shocks affecting discounting from shocks changing cash flows.
\item \textbf{Binding and segmented markets (editorial).} Permit changing marginal intermediaries and entry; identify finite or one-sided return responses at leverage boundaries.
\end{enumerate}
\subsection*{References and source audit}
\begin{enumerate}
\item Official January 9, 2026 web version; locator: The Financial System / Financial market dynamics. Broad agenda; exact model editorial. URL: \url{https://www.bankofengland.co.uk/research/bank-of-england-agenda-for-research}.
\item Publisher metadata and abstract verified. Locator: abstract. Background only; no exact open passage verified. URL: \url{https://www.annualreviews.org/content/journals/10.1146/annurev-financial-090524-120754}.
\end{enumerate}
\begin{enumerate}\raggedright
\item\hypertarget{definitions:source:30007598:1}{} Bank of England. 2026. \emph{Bank of England Agenda for Research, 2025–2028}. The Financial System / Financial market dynamics.
\url{https://www.bankofengland.co.uk/research/bank-of-england-agenda-for-research}.
\item\hypertarget{definitions:source:30007598:2}{} Valentin Haddad; Tyler Muir. 2025. \emph{Market Macrostructure: Institutions and Asset Prices}. Publisher abstract.
\url{https://www.annualreviews.org/content/journals/10.1146/annurev-financial-090524-120754}.
DOI: \url{https://doi.org/10.1146/annurev-financial-090524-120754}.
\end{enumerate}

\subsection*{Common conventions: Source mathematical conventions}

These conventions apply to each entry unless explicitly replaced.

\textbf{Models and identification.} A primitive model \(m\) specifies all probability laws, feasible choices, information, equilibrium and measurement rules used by its target. The admissible class \(\mathcal M\) is fixed independently of outcomes. If \(P_O(m)\) is its observable probability law and \(\tau(m)\) the target, its identified set at observable law \(P\) is
\[
 \mathcal I_\tau(P)=\{\tau(m):m\in\mathcal M,\ P_O(m)=P\}.
\]
Point identification means that this set is a singleton. An empty set signals incompatibility of data and assumptions. Coordinate bounds need not describe the joint set. Bounds are sharp only if every retained target is attainable or arbitrarily approachable by compatible models. Finite bounds require explicit restrictions on unobserved quantities; unrestricted confounding, hidden wealth, extrapolation or arbitrary equilibrium selection can yield uninformative sets.

\textbf{Causal interventions.} A potential outcome \(Y_i(p)\) is the outcome under a complete policy regime \(p\), including timing, financing, information and market spillovers. The target population and units are fixed before treatment. With interference, use the complete assignment vector or a justified exposure mapping; individual no-interference notation alone cannot identify an economy-wide intervention. A difference of mean potential outcomes does not require the cross-world joint distribution. Conditioning on post-treatment migration, survival, enrollment or employment changes the estimand and needs separate justification.

\textbf{Statistical statements.} An estimator, confidence set, forecast or test specifies a sampling law, model class and loss/coverage criterion. Uniform \((1-\alpha)\)-coverage means \(\inf_{m\in\mathcal M}\Pr_m\{\tau(m)\in C_n\}\geq1-\alpha\), or an explicitly asymptotic version. The whole parameter space trivially has coverage, so informativeness requires a stated width, risk or power objective. A forecasting improvement is relative to a named benchmark, information set and evaluation population.

\textbf{Equilibrium and optimization.} An equilibrium comprises feasible plans, optimality, beliefs where needed, budget constraints and all market/resource clearing. The primitives or a stated benchmark define each correspondence \(\mathcal E(p;m)\). Nonemptiness is assumed or proved, never inferred just from compact policy sets. A selected-equilibrium target uses a declared selection rule; a robust target quantifies over all permitted equilibria. Use suprema or infima unless attainment is established. An \(\varepsilon\)-optimal policy is feasible and within \(\varepsilon\) of the finite optimum value. Report an empty feasible set separately.

\textbf{Welfare and accounting.} Normative utility, interpersonal normalization, social weights, numeraire, discounting and the use of public receipts are primitives. Behavioral utility need not coincide with the welfare criterion. Household welfare includes ownership income and rents once; adding profits again to compensating variation generally double-counts them. A monetary equivalent is defined by an explicit transfer and price convention, with monotonicity and range conditions. Average effects, mechanism contributions, transfers and welfare losses are different targets. Infinite sums require convergence and dynamic feasible plans require terminal or no-Ponzi conditions.

\textbf{Source versions.} A working-paper date, revision date and journal publication year are distinguished. A DOI for a journal version is not automatically the DOI of an earlier manuscript. Locators refer to the stated version and printed pages where specified. A metadata-only check does not verify a claim about a full-text conclusion. Every entry's source subsection narrows the provenance claim accordingly.
% END ECONOMICS STATEMENT
\end{document}
